\documentclass[11pt,a4paper]{article}
\newcommand{\gradelabel}{Grade 10}
\newcommand{\leconhdr}{Variables et types de données}
% Préambule commun aux plans de cours de robotique (Grades 7 à 12).
% Chaque plan de cours commence par \documentclass[11pt,a4paper]{article}, définit
% \gradelabel et \leconhdr (en-têtes de page), puis % Préambule commun aux plans de cours de robotique (Grades 7 à 12).
% Chaque plan de cours commence par \documentclass[11pt,a4paper]{article}, définit
% \gradelabel et \leconhdr (en-têtes de page), puis % Préambule commun aux plans de cours de robotique (Grades 7 à 12).
% Chaque plan de cours commence par \documentclass[11pt,a4paper]{article}, définit
% \gradelabel et \leconhdr (en-têtes de page), puis \input{../../preambule_robotique.tex}.
\usepackage[T1]{fontenc}
\usepackage[utf8]{inputenc}
\usepackage{lmodern}
\usepackage[french]{babel}
\usepackage[margin=2cm,headheight=14pt]{geometry}
\usepackage[most]{tcolorbox}
\usepackage{xcolor,tabularx,booktabs,array,enumitem,fancyhdr,amssymb,tikz}
\usetikzlibrary{arrows.meta,positioning}

\pagestyle{fancy}
\fancyhf{}
\lhead{\small Robotique -- \gradelabel{} -- Novembre 2026}
\rhead{\small \leconhdr}
\cfoot{\thepage}
\setlist{itemsep=1pt,topsep=2pt,leftmargin=*}
\renewcommand{\arraystretch}{1.35}
\newcolumntype{Y}{>{\raggedright\arraybackslash}X}
\newcommand{\cb}{$\square$\,}
\newcommand{\fl}{$\rightarrow$}
\newcommand{\lig}{\noindent\rule{\linewidth}{0.3pt}\\[5mm]}

\newtcolorbox{etape}[3]{enhanced,breakable,colback=#1!6,colframe=#1!70!black,
  coltitle=white,fonttitle=\bfseries,title={#2\hfill #3},boxrule=0.8pt,arc=2mm,
  before skip=8pt,after skip=8pt}
\newtcolorbox{astuce}[1][Différenciation et besoins particuliers]{breakable,
  colback=yellow!12,colframe=yellow!50!black,fonttitle=\bfseries\small,
  title={#1},boxrule=0.5pt,arc=1mm,before skip=6pt,after skip=6pt}
\newtcolorbox{code}[1][Code]{breakable,colback=black!4,colframe=black!60,
  fonttitle=\bfseries\small,title={#1},boxrule=0.5pt,arc=1mm,fontupper=\ttfamily\small,
  before skip=6pt,after skip=6pt}
% Blocs de code coloré (écrire le code sans accents : comme dans l'éditeur)
\newtcblisting{arduino}[1][Code Arduino]{listing only,breakable,colback=black!4,
  colframe=black!60,fonttitle=\bfseries\small,title={#1},boxrule=0.5pt,arc=1mm,
  before skip=6pt,after skip=6pt,
  listing options={language=C++,basicstyle=\ttfamily\small,columns=fullflexible,
  keepspaces=true,showstringspaces=false,commentstyle=\color{green!40!black}\itshape,
  keywordstyle=\color{blue!70!black}\bfseries}}
\newtcblisting{python}[1][Code Python]{listing only,breakable,colback=black!4,
  colframe=black!60,fonttitle=\bfseries\small,title={#1},boxrule=0.5pt,arc=1mm,
  before skip=6pt,after skip=6pt,
  listing options={language=Python,basicstyle=\ttfamily\small,columns=fullflexible,
  keepspaces=true,showstringspaces=false,commentstyle=\color{green!40!black}\itshape,
  keywordstyle=\color{blue!70!black}\bfseries}}
\newcommand{\entete}[2]{%
  \begin{tcolorbox}[colback=#1!10,colframe=#1!70!black,boxrule=0.8pt,arc=2mm]
  {\Large\bfseries #2}\hfill Nom : \rule{4cm}{0.3pt}\quad Date : \rule{2.2cm}{0.3pt}
  \end{tcolorbox}}

% Bandeau de titre : #1 titre, #2 semaine
\newcommand{\titreblock}[2]{%
\begin{tcolorbox}[colback=teal!10,colframe=teal!70!black,arc=2mm,boxrule=1pt]
\centering
{\LARGE\bfseries #1}\\[2mm]
{\large Robotique -- \gradelabel{} -- Plan de cours de 2 heures}\\[1mm]
Enseignant : Francis Ki \quad|\quad #2
\end{tcolorbox}}

% Tableau du minutage : 6 contenus (accroche, mini-leçon, guidée, autonome, mise en commun, exit ticket)
\newcommand{\minutage}[6]{%
\begin{tabularx}{\linewidth}{>{\bfseries}p{2.2cm}p{2.6cm}Y p{3cm}}
\toprule
Minutes & Étape & Contenu & Organisation\\
\midrule
0 -- 10 & Accroche & #1 & Classe entière\\
10 -- 30 & Mini-leçon & #2 & Classe entière\\
30 -- 60 & Pratique guidée & #3 & Collectif puis binômes\\
60 -- 95 & Pratique autonome & #4 & Individuel sur poste\\
95 -- 110 & Mise en commun & #5 & Classe entière\\
110 -- 120 & Exit ticket & #6 & Individuel sur poste\\
\bottomrule
\end{tabularx}}

% Étayage progressif : #1 aide 1, #2 aide 2, #3 aide 3
\newcommand{\aides}[3]{%
\textbf{Étayage progressif de l'aide.}
\begin{center}
\begin{tabularx}{\linewidth}{>{\bfseries}p{2.4cm}Y p{4.4cm}}
\toprule
Niveau d'aide & Contenu & Qui intervient\\
\midrule
Aide 1 & #1 & Le binôme seul\\
Aide 2 & #2 & Le binôme avec la carte d'aide\\
Aide 3 & #3 & L'enseignant\\
\bottomrule
\end{tabularx}
\end{center}}

% Les trois groupes de la pratique autonome : #1 tâche A, #2 tâche B, #3 tâche C
\newcommand{\groupes}[3]{%
\textbf{Consigne commune\,:} chaque élève ouvre le fichier de son groupe, le renomme avec son prénom et travaille seul sur son poste. L'enseignant passe d'abord auprès du groupe A pendant 10 min.

\medskip
\begin{tabularx}{\linewidth}{>{\bfseries}p{3cm}Y p{2.2cm}}
\toprule
Groupe & Tâche sur poste & Annexe\\
\midrule
A : en dessous du niveau & #1 & Annexe C\\
B : au niveau & #2 & Annexe D\\
C : au-dessus du niveau & #3 & Annexe E\\
\bottomrule
\end{tabularx}}

% Mise en commun : #1 groupe A, #2 groupe B, #3 groupe C, #4 conclusion
\newcommand{\miseencommun}[4]{%
\begin{etape}{teal}{5. Mise en commun}{95 -- 110 min}
\begin{enumerate}
 \item (2 min) Retour au calme, écrans mis en veille.
 \item (9 min) Trois passages courts de 3 min, avec valorisation :
 \begin{itemize}
  \item \textbf{Groupe A} : #1
  \item \textbf{Groupe B} : #2
  \item \textbf{Groupe C} : #3
 \end{itemize}
 \item (4 min) Conclusion : #4
\end{enumerate}
\end{etape}}

% Clôture : #1 question 1, #2 question 2
\newcommand{\cloture}[2]{%
\begin{etape}{gray}{6. Clôture : exit ticket}{110 -- 120 min}
Formulaire en ligne ou papier (annexe G) :
\begin{enumerate}
 \item #1
 \item #2
 \item Auto-évaluation : \emph{Je maîtrise / Je progresse / J'ai besoin d'aide.}
\end{enumerate}
Lire les réponses avant la séance suivante pour ajuster les groupes de besoin.
\end{etape}}

% Grille d'observation : #1 #2 #3 intitulés de colonnes
\newcommand{\grille}[3]{%
\subsection*{Grille d'observation (enseignant)}
\begin{tabularx}{\linewidth}{Y p{2.4cm} p{2.4cm} p{2.4cm} Y}
\toprule
Élève & #1 & #2 & #3 & Groupe suivant\\
\midrule
 & & & & \\ & & & & \\ & & & & \\ & & & & \\ & & & & \\
\bottomrule
\end{tabularx}}

% Signature de l'exit ticket
\newcommand{\autoeval}{Je me sens : \cb Je maîtrise \quad \cb Je progresse \quad \cb J'ai besoin d'aide}
.
\usepackage[T1]{fontenc}
\usepackage[utf8]{inputenc}
\usepackage{lmodern}
\usepackage[french]{babel}
\usepackage[margin=2cm,headheight=14pt]{geometry}
\usepackage[most]{tcolorbox}
\usepackage{xcolor,tabularx,booktabs,array,enumitem,fancyhdr,amssymb,tikz}
\usetikzlibrary{arrows.meta,positioning}

\pagestyle{fancy}
\fancyhf{}
\lhead{\small Robotique -- \gradelabel{} -- Novembre 2026}
\rhead{\small \leconhdr}
\cfoot{\thepage}
\setlist{itemsep=1pt,topsep=2pt,leftmargin=*}
\renewcommand{\arraystretch}{1.35}
\newcolumntype{Y}{>{\raggedright\arraybackslash}X}
\newcommand{\cb}{$\square$\,}
\newcommand{\fl}{$\rightarrow$}
\newcommand{\lig}{\noindent\rule{\linewidth}{0.3pt}\\[5mm]}

\newtcolorbox{etape}[3]{enhanced,breakable,colback=#1!6,colframe=#1!70!black,
  coltitle=white,fonttitle=\bfseries,title={#2\hfill #3},boxrule=0.8pt,arc=2mm,
  before skip=8pt,after skip=8pt}
\newtcolorbox{astuce}[1][Différenciation et besoins particuliers]{breakable,
  colback=yellow!12,colframe=yellow!50!black,fonttitle=\bfseries\small,
  title={#1},boxrule=0.5pt,arc=1mm,before skip=6pt,after skip=6pt}
\newtcolorbox{code}[1][Code]{breakable,colback=black!4,colframe=black!60,
  fonttitle=\bfseries\small,title={#1},boxrule=0.5pt,arc=1mm,fontupper=\ttfamily\small,
  before skip=6pt,after skip=6pt}
% Blocs de code coloré (écrire le code sans accents : comme dans l'éditeur)
\newtcblisting{arduino}[1][Code Arduino]{listing only,breakable,colback=black!4,
  colframe=black!60,fonttitle=\bfseries\small,title={#1},boxrule=0.5pt,arc=1mm,
  before skip=6pt,after skip=6pt,
  listing options={language=C++,basicstyle=\ttfamily\small,columns=fullflexible,
  keepspaces=true,showstringspaces=false,commentstyle=\color{green!40!black}\itshape,
  keywordstyle=\color{blue!70!black}\bfseries}}
\newtcblisting{python}[1][Code Python]{listing only,breakable,colback=black!4,
  colframe=black!60,fonttitle=\bfseries\small,title={#1},boxrule=0.5pt,arc=1mm,
  before skip=6pt,after skip=6pt,
  listing options={language=Python,basicstyle=\ttfamily\small,columns=fullflexible,
  keepspaces=true,showstringspaces=false,commentstyle=\color{green!40!black}\itshape,
  keywordstyle=\color{blue!70!black}\bfseries}}
\newcommand{\entete}[2]{%
  \begin{tcolorbox}[colback=#1!10,colframe=#1!70!black,boxrule=0.8pt,arc=2mm]
  {\Large\bfseries #2}\hfill Nom : \rule{4cm}{0.3pt}\quad Date : \rule{2.2cm}{0.3pt}
  \end{tcolorbox}}

% Bandeau de titre : #1 titre, #2 semaine
\newcommand{\titreblock}[2]{%
\begin{tcolorbox}[colback=teal!10,colframe=teal!70!black,arc=2mm,boxrule=1pt]
\centering
{\LARGE\bfseries #1}\\[2mm]
{\large Robotique -- \gradelabel{} -- Plan de cours de 2 heures}\\[1mm]
Enseignant : Francis Ki \quad|\quad #2
\end{tcolorbox}}

% Tableau du minutage : 6 contenus (accroche, mini-leçon, guidée, autonome, mise en commun, exit ticket)
\newcommand{\minutage}[6]{%
\begin{tabularx}{\linewidth}{>{\bfseries}p{2.2cm}p{2.6cm}Y p{3cm}}
\toprule
Minutes & Étape & Contenu & Organisation\\
\midrule
0 -- 10 & Accroche & #1 & Classe entière\\
10 -- 30 & Mini-leçon & #2 & Classe entière\\
30 -- 60 & Pratique guidée & #3 & Collectif puis binômes\\
60 -- 95 & Pratique autonome & #4 & Individuel sur poste\\
95 -- 110 & Mise en commun & #5 & Classe entière\\
110 -- 120 & Exit ticket & #6 & Individuel sur poste\\
\bottomrule
\end{tabularx}}

% Étayage progressif : #1 aide 1, #2 aide 2, #3 aide 3
\newcommand{\aides}[3]{%
\textbf{Étayage progressif de l'aide.}
\begin{center}
\begin{tabularx}{\linewidth}{>{\bfseries}p{2.4cm}Y p{4.4cm}}
\toprule
Niveau d'aide & Contenu & Qui intervient\\
\midrule
Aide 1 & #1 & Le binôme seul\\
Aide 2 & #2 & Le binôme avec la carte d'aide\\
Aide 3 & #3 & L'enseignant\\
\bottomrule
\end{tabularx}
\end{center}}

% Les trois groupes de la pratique autonome : #1 tâche A, #2 tâche B, #3 tâche C
\newcommand{\groupes}[3]{%
\textbf{Consigne commune\,:} chaque élève ouvre le fichier de son groupe, le renomme avec son prénom et travaille seul sur son poste. L'enseignant passe d'abord auprès du groupe A pendant 10 min.

\medskip
\begin{tabularx}{\linewidth}{>{\bfseries}p{3cm}Y p{2.2cm}}
\toprule
Groupe & Tâche sur poste & Annexe\\
\midrule
A : en dessous du niveau & #1 & Annexe C\\
B : au niveau & #2 & Annexe D\\
C : au-dessus du niveau & #3 & Annexe E\\
\bottomrule
\end{tabularx}}

% Mise en commun : #1 groupe A, #2 groupe B, #3 groupe C, #4 conclusion
\newcommand{\miseencommun}[4]{%
\begin{etape}{teal}{5. Mise en commun}{95 -- 110 min}
\begin{enumerate}
 \item (2 min) Retour au calme, écrans mis en veille.
 \item (9 min) Trois passages courts de 3 min, avec valorisation :
 \begin{itemize}
  \item \textbf{Groupe A} : #1
  \item \textbf{Groupe B} : #2
  \item \textbf{Groupe C} : #3
 \end{itemize}
 \item (4 min) Conclusion : #4
\end{enumerate}
\end{etape}}

% Clôture : #1 question 1, #2 question 2
\newcommand{\cloture}[2]{%
\begin{etape}{gray}{6. Clôture : exit ticket}{110 -- 120 min}
Formulaire en ligne ou papier (annexe G) :
\begin{enumerate}
 \item #1
 \item #2
 \item Auto-évaluation : \emph{Je maîtrise / Je progresse / J'ai besoin d'aide.}
\end{enumerate}
Lire les réponses avant la séance suivante pour ajuster les groupes de besoin.
\end{etape}}

% Grille d'observation : #1 #2 #3 intitulés de colonnes
\newcommand{\grille}[3]{%
\subsection*{Grille d'observation (enseignant)}
\begin{tabularx}{\linewidth}{Y p{2.4cm} p{2.4cm} p{2.4cm} Y}
\toprule
Élève & #1 & #2 & #3 & Groupe suivant\\
\midrule
 & & & & \\ & & & & \\ & & & & \\ & & & & \\ & & & & \\
\bottomrule
\end{tabularx}}

% Signature de l'exit ticket
\newcommand{\autoeval}{Je me sens : \cb Je maîtrise \quad \cb Je progresse \quad \cb J'ai besoin d'aide}
.
\usepackage[T1]{fontenc}
\usepackage[utf8]{inputenc}
\usepackage{lmodern}
\usepackage[french]{babel}
\usepackage[margin=2cm,headheight=14pt]{geometry}
\usepackage[most]{tcolorbox}
\usepackage{xcolor,tabularx,booktabs,array,enumitem,fancyhdr,amssymb,tikz}
\usetikzlibrary{arrows.meta,positioning}

\pagestyle{fancy}
\fancyhf{}
\lhead{\small Robotique -- \gradelabel{} -- Novembre 2026}
\rhead{\small \leconhdr}
\cfoot{\thepage}
\setlist{itemsep=1pt,topsep=2pt,leftmargin=*}
\renewcommand{\arraystretch}{1.35}
\newcolumntype{Y}{>{\raggedright\arraybackslash}X}
\newcommand{\cb}{$\square$\,}
\newcommand{\fl}{$\rightarrow$}
\newcommand{\lig}{\noindent\rule{\linewidth}{0.3pt}\\[5mm]}

\newtcolorbox{etape}[3]{enhanced,breakable,colback=#1!6,colframe=#1!70!black,
  coltitle=white,fonttitle=\bfseries,title={#2\hfill #3},boxrule=0.8pt,arc=2mm,
  before skip=8pt,after skip=8pt}
\newtcolorbox{astuce}[1][Différenciation et besoins particuliers]{breakable,
  colback=yellow!12,colframe=yellow!50!black,fonttitle=\bfseries\small,
  title={#1},boxrule=0.5pt,arc=1mm,before skip=6pt,after skip=6pt}
\newtcolorbox{code}[1][Code]{breakable,colback=black!4,colframe=black!60,
  fonttitle=\bfseries\small,title={#1},boxrule=0.5pt,arc=1mm,fontupper=\ttfamily\small,
  before skip=6pt,after skip=6pt}
% Blocs de code coloré (écrire le code sans accents : comme dans l'éditeur)
\newtcblisting{arduino}[1][Code Arduino]{listing only,breakable,colback=black!4,
  colframe=black!60,fonttitle=\bfseries\small,title={#1},boxrule=0.5pt,arc=1mm,
  before skip=6pt,after skip=6pt,
  listing options={language=C++,basicstyle=\ttfamily\small,columns=fullflexible,
  keepspaces=true,showstringspaces=false,commentstyle=\color{green!40!black}\itshape,
  keywordstyle=\color{blue!70!black}\bfseries}}
\newtcblisting{python}[1][Code Python]{listing only,breakable,colback=black!4,
  colframe=black!60,fonttitle=\bfseries\small,title={#1},boxrule=0.5pt,arc=1mm,
  before skip=6pt,after skip=6pt,
  listing options={language=Python,basicstyle=\ttfamily\small,columns=fullflexible,
  keepspaces=true,showstringspaces=false,commentstyle=\color{green!40!black}\itshape,
  keywordstyle=\color{blue!70!black}\bfseries}}
\newcommand{\entete}[2]{%
  \begin{tcolorbox}[colback=#1!10,colframe=#1!70!black,boxrule=0.8pt,arc=2mm]
  {\Large\bfseries #2}\hfill Nom : \rule{4cm}{0.3pt}\quad Date : \rule{2.2cm}{0.3pt}
  \end{tcolorbox}}

% Bandeau de titre : #1 titre, #2 semaine
\newcommand{\titreblock}[2]{%
\begin{tcolorbox}[colback=teal!10,colframe=teal!70!black,arc=2mm,boxrule=1pt]
\centering
{\LARGE\bfseries #1}\\[2mm]
{\large Robotique -- \gradelabel{} -- Plan de cours de 2 heures}\\[1mm]
Enseignant : Francis Ki \quad|\quad #2
\end{tcolorbox}}

% Tableau du minutage : 6 contenus (accroche, mini-leçon, guidée, autonome, mise en commun, exit ticket)
\newcommand{\minutage}[6]{%
\begin{tabularx}{\linewidth}{>{\bfseries}p{2.2cm}p{2.6cm}Y p{3cm}}
\toprule
Minutes & Étape & Contenu & Organisation\\
\midrule
0 -- 10 & Accroche & #1 & Classe entière\\
10 -- 30 & Mini-leçon & #2 & Classe entière\\
30 -- 60 & Pratique guidée & #3 & Collectif puis binômes\\
60 -- 95 & Pratique autonome & #4 & Individuel sur poste\\
95 -- 110 & Mise en commun & #5 & Classe entière\\
110 -- 120 & Exit ticket & #6 & Individuel sur poste\\
\bottomrule
\end{tabularx}}

% Étayage progressif : #1 aide 1, #2 aide 2, #3 aide 3
\newcommand{\aides}[3]{%
\textbf{Étayage progressif de l'aide.}
\begin{center}
\begin{tabularx}{\linewidth}{>{\bfseries}p{2.4cm}Y p{4.4cm}}
\toprule
Niveau d'aide & Contenu & Qui intervient\\
\midrule
Aide 1 & #1 & Le binôme seul\\
Aide 2 & #2 & Le binôme avec la carte d'aide\\
Aide 3 & #3 & L'enseignant\\
\bottomrule
\end{tabularx}
\end{center}}

% Les trois groupes de la pratique autonome : #1 tâche A, #2 tâche B, #3 tâche C
\newcommand{\groupes}[3]{%
\textbf{Consigne commune\,:} chaque élève ouvre le fichier de son groupe, le renomme avec son prénom et travaille seul sur son poste. L'enseignant passe d'abord auprès du groupe A pendant 10 min.

\medskip
\begin{tabularx}{\linewidth}{>{\bfseries}p{3cm}Y p{2.2cm}}
\toprule
Groupe & Tâche sur poste & Annexe\\
\midrule
A : en dessous du niveau & #1 & Annexe C\\
B : au niveau & #2 & Annexe D\\
C : au-dessus du niveau & #3 & Annexe E\\
\bottomrule
\end{tabularx}}

% Mise en commun : #1 groupe A, #2 groupe B, #3 groupe C, #4 conclusion
\newcommand{\miseencommun}[4]{%
\begin{etape}{teal}{5. Mise en commun}{95 -- 110 min}
\begin{enumerate}
 \item (2 min) Retour au calme, écrans mis en veille.
 \item (9 min) Trois passages courts de 3 min, avec valorisation :
 \begin{itemize}
  \item \textbf{Groupe A} : #1
  \item \textbf{Groupe B} : #2
  \item \textbf{Groupe C} : #3
 \end{itemize}
 \item (4 min) Conclusion : #4
\end{enumerate}
\end{etape}}

% Clôture : #1 question 1, #2 question 2
\newcommand{\cloture}[2]{%
\begin{etape}{gray}{6. Clôture : exit ticket}{110 -- 120 min}
Formulaire en ligne ou papier (annexe G) :
\begin{enumerate}
 \item #1
 \item #2
 \item Auto-évaluation : \emph{Je maîtrise / Je progresse / J'ai besoin d'aide.}
\end{enumerate}
Lire les réponses avant la séance suivante pour ajuster les groupes de besoin.
\end{etape}}

% Grille d'observation : #1 #2 #3 intitulés de colonnes
\newcommand{\grille}[3]{%
\subsection*{Grille d'observation (enseignant)}
\begin{tabularx}{\linewidth}{Y p{2.4cm} p{2.4cm} p{2.4cm} Y}
\toprule
Élève & #1 & #2 & #3 & Groupe suivant\\
\midrule
 & & & & \\ & & & & \\ & & & & \\ & & & & \\ & & & & \\
\bottomrule
\end{tabularx}}

% Signature de l'exit ticket
\newcommand{\autoeval}{Je me sens : \cb Je maîtrise \quad \cb Je progresse \quad \cb J'ai besoin d'aide}


\begin{document}
\titreblock{Variables et types de données (Python)}{Semaine du 2 au 6 novembre 2026}

%==================================================================
\section{Informations générales}
\begin{tabularx}{\linewidth}{>{\bfseries}p{3.6cm}Y}
\toprule
Titre & Variables et types de données en Python\\
Classe / effectif & Grade 10, 11 à 20 élèves, niveaux hétérogènes, élèves à besoins particuliers\\
Durée & 2 heures (120 min)\\
Lieu & Salle informatique, un poste par élève\\
Matériel et logiciels &
\begin{itemize}
 \item Un poste par élève avec l'éditeur \textbf{Mu} en mode \emph{Python 3} (boutons Exécuter, REPL, Vérifier, Débogueur).
 \item Vidéoprojecteur ou tableau interactif avec Mu projeté.
 \item Un dossier partagé contenant les fichiers de départ \texttt{.py} des trois niveaux.
 \item Inspiration pédagogique : unité «~Introduction to Python programming~» de TeachComputing (programmation en binômes, codage en direct, exemples commentés, traitement des idées fausses courantes).
\end{itemize}\\
Pré-requis &
Séance d'octobre : installer et lancer Mu, exécuter un programme simple avec \texttt{print}. Notions de réseau vues en septembre non nécessaires.\\
\bottomrule
\end{tabularx}

\begin{astuce}[Préparation avant le cours (environ 40 min)]
\begin{enumerate}
 \item Vérifier que Mu démarre en mode Python 3 sur tous les postes et que l'exécution d'un \texttt{print} fonctionne.
 \item Préparer les fichiers de départ : \texttt{niveauA.py} (programme à compléter), \texttt{niveauB.py} (vide avec commentaires), \texttt{niveauC.py} (programme contenant des bogues volontaires).
 \item Créer le formulaire de sortie (annexe G) ou imprimer l'annexe.
 \item Préparer les diapositives de codage en direct (trois exemples courts).
\end{enumerate}
\end{astuce}

%==================================================================
\section{Objectifs d'apprentissage}
\begin{tcolorbox}[colback=gray!8,colframe=gray!60!black,arc=1mm,boxrule=0.5pt]
\textbf{Objectifs du programme}
\begin{itemize}
 \item Déclarer des variables.
 \item Manipuler des nombres et du texte.
\end{itemize}
\textbf{À la fin de la séance, l'élève est capable\,:}
\begin{itemize}
 \item de créer une variable avec l'affectation \texttt{=} et d'afficher sa valeur\,;
 \item de distinguer les types \texttt{int}, \texttt{float}, \texttt{str} et \texttt{bool}, et de convertir une saisie avec \texttt{int()} ou \texttt{float()}\,;
 \item de combiner texte et nombres avec les f-strings et de corriger l'erreur classique «~\texttt{'5' + '3'}~».
\end{itemize}
\textbf{Compétences numériques\,:} saisir et exécuter du code dans Mu, lire un message d'erreur, nommer ses fichiers.
\end{tcolorbox}

%==================================================================
\section{Déroulé minuté (120 min)}
\minutage{Prédire avant d'exécuter : 5 + 3 ou \texttt{'5' + '3'} ?}{Variable, type, saisie et f-string}{Un programme «~fiche d'identité~» en collectif, puis un calcul d'aire en binômes}{Trois niveaux de besoin}{Valoriser chaque groupe}{Trois questions sur poste}

%------------------------------------------------------------------
\begin{etape}{orange}{1. Accroche : prédire avant d'exécuter}{0 -- 10 min}
\begin{enumerate}
 \item (5 min) Projeter deux lignes : \texttt{print(5 + 3)} et \texttt{print('5' + '3')}. Chaque élève écrit sa prédiction, puis lève la main pour 8 ou 53.
 \item (5 min) Exécuter en direct. Question-problème : \textbf{«~Pourquoi l'ordinateur répond-il différemment à deux calculs qui se ressemblent ?~»}
\end{enumerate}
\end{etape}

%------------------------------------------------------------------
\begin{etape}{blue}{2. Mini-leçon : variables et types}{10 -- 30 min}
\textbf{Support\,:} annexe A. Codage en direct de l'enseignant. Quatre temps de 5 min.

\textbf{Temps 1 : une variable est une boîte étiquetée.}
\begin{python}[Affectation]
age = 15            # on range 15 dans la boite age
prenom = "Awa"      # on range du texte
print(prenom, age)
\end{python}
Le signe \texttt{=} n'est pas l'égalité des mathématiques : il \emph{range} la valeur de droite dans la variable de gauche.

\textbf{Temps 2 : les types.}
\begin{center}
\begin{tabularx}{\linewidth}{>{\ttfamily}p{1.8cm}Y Y}
\toprule
Type & Exemple & Usage\\
\midrule
int & 15 & Nombre entier\\
float & 3.14 & Nombre à virgule (point décimal)\\
str & "Awa" & Texte (entre guillemets)\\
bool & True & Vrai ou faux\\
\bottomrule
\end{tabularx}
\end{center}

\textbf{Temps 3 : \texttt{input} renvoie toujours du texte.}
\begin{python}[Conversion]
a = input("Premier nombre : ")   # a est un str
b = input("Second nombre : ")
somme = int(a) + int(b)          # conversion en entier
print(f"La somme est {somme}")
\end{python}

\textbf{Temps 4 : opérations utiles.} \texttt{+ - * /} ; \texttt{//} (division entière) ; \texttt{\%} (reste) ; \texttt{**} (puissance). Le type d'une variable s'obtient avec \texttt{type(x)}.

\begin{astuce}
Garder l'annexe A à côté du clavier. Insister sur : un nom de variable sans espace ni accent, majuscules et minuscules différentes. Lire le code à voix haute en disant «~range dans~» pour \texttt{=}. Autoriser le travail avec un camarade comme pilote.
\end{astuce}
\end{etape}

%------------------------------------------------------------------
\begin{etape}{green}{3. Pratique guidée}{30 -- 60 min}
\textbf{Fiche\,:} annexe B. Programmation en binômes (pilote et navigateur, rôles échangés à mi-temps).

\textbf{3a. Fiche d'identité en collectif (12 min).} Écrire ensemble, en codage en direct, un programme qui demande le prénom et l'âge puis affiche une phrase avec une f-string et l'âge dans cinq ans.

\textbf{3b. Aire d'un rectangle en binômes (18 min).} Demander la longueur et la largeur (avec \texttt{float}), calculer l'aire et le périmètre, puis les afficher avec une phrase.

\aides{Le programme du 3a affiché, l'annexe A et la liste des types}{Questions guides : \emph{Quel est le type de la saisie ? Que faut-il faire avant de calculer ? Quelle valeur la variable contient-elle ?}}{L'enseignant tape la première ligne avec le binôme, qui continue}

\begin{astuce}
Faire expliquer à voix haute chaque ligne par le navigateur pendant que le pilote tape. Rappeler le bouton Vérifier de Mu pour repérer les fautes de syntaxe avant d'exécuter.
\end{astuce}
\end{etape}

%------------------------------------------------------------------
\begin{etape}{violet}{4. Pratique autonome différenciée}{60 -- 95 min}
\groupes{Fichier \texttt{niveauA.py} à compléter : lignes à trous avec une banque de mots (\texttt{input}, \texttt{int}, \texttt{print}, \texttt{prenom}, \texttt{age}), puis appariement valeur et type. Exécuter et comparer avec le résultat attendu.}{Écrire seul un programme qui calcule la moyenne de trois notes saisies et affiche une phrase avec la moyenne arrondie à une décimale. Tester avec trois jeux de valeurs.}{Corriger les bogues du fichier \texttt{niveauC.py} (types oubliés, nom de variable mal orthographié, division mal écrite) et écrire un convertisseur de température Celsius vers Fahrenheit avec affichage formaté. Expliquer chaque correction en commentaire.}

\medskip
\textbf{Plan B si un élève termine tôt\,:} calculer l'IMC (masse divisée par taille au carré) et l'afficher avec deux décimales.
\end{etape}

%------------------------------------------------------------------
\miseencommun{un élève lit son programme complété et explique chaque ligne. Valoriser la lecture à voix haute du code.}{un binôme montre sa moyenne et ses trois tests. Valoriser les tests variés.}{un élève présente un bogue corrigé et explique ce qui n'allait pas. Valoriser la méthode de recherche d'erreur.}{\textit{une variable range une valeur, chaque valeur a un type, et \texttt{input} donne du texte qu'il faut convertir pour calculer.}}

\cloture{Quel est le type de la variable \texttt{x} après \texttt{x = input("Nombre : ")} ?}{Que vaut \texttt{int("4") + 3} ? Et \texttt{"4" + "3"} ?}

%==================================================================
% ANNEXES
%==================================================================
\newpage
\section*{Annexe A : Support résumé de la mini-leçon}
\entete{blue}{Variables et types en Python}

\begin{tcolorbox}[colback=yellow!10,colframe=yellow!50!black]
\texttt{=} range une valeur dans une variable. Chaque valeur a un \textbf{type}. \texttt{input()} renvoie du \textbf{texte} : il faut convertir avec \texttt{int()} ou \texttt{float()}.
\end{tcolorbox}

\begin{tabularx}{\linewidth}{>{\ttfamily}p{1.8cm}Y Y}
\toprule
Type & Exemple & Convertir\\
\midrule
int & 15 & int("15")\\
float & 3.14 & float("3.14")\\
str & "Awa" & str(15)\\
bool & True & bool(1)\\
\bottomrule
\end{tabularx}

\begin{python}[Modèle]
prenom = input("Ton prenom : ")
age = int(input("Ton age : "))
print(f"{prenom} aura {age + 5} ans dans 5 ans.")
\end{python}

\begin{tabularx}{\linewidth}{>{\ttfamily}p{2.2cm}Y}
\toprule
Opérateur & Rôle\\
\midrule
+ - * / & Opérations usuelles\\
// & Division entière\\
\% & Reste de la division\\
** & Puissance\\
\bottomrule
\end{tabularx}

%==================================================================
\newpage
\section*{Annexe B : Fiche de pratique guidée}
\entete{green}{Pratique guidée : mon premier programme}

\subsection*{3a. Fiche d'identité (collectif)}
\begin{python}[À compléter pendant le codage en direct]
prenom = input("Ton prenom : ")
age = _______(input("Ton age : "))
print(f"Bonjour {_______}, tu auras {age + 5} ans dans 5 ans.")
\end{python}

\subsection*{3b. Aire d'un rectangle (binôme)}
\begin{tabularx}{\linewidth}{>{\bfseries}p{3.6cm}Y Y}
\toprule
Étape & Mon code & Type de la variable\\
\midrule
Demander la longueur & & \\
Demander la largeur & & \\
Calculer l'aire & & \\
Calculer le périmètre & & \\
Afficher la phrase & & \\
\bottomrule
\end{tabularx}

\begin{tcolorbox}[colback=yellow!10,colframe=yellow!50!black,title={Cartes d'aide à utiliser dans l'ordre},fonttitle=\bfseries\small]
\textbf{Aide 1 :} regarde le programme du 3a et l'annexe A.\\
\textbf{Aide 2 :} Quel est le type de la saisie ? Que faut-il faire avant de calculer ?\\
\textbf{Aide 3 :} demande à l'enseignant de taper la première ligne avec toi.
\end{tcolorbox}

%==================================================================
\newpage
\section*{Annexe C : Niveau A (en dessous du niveau)}
\entete{orange}{Niveau A : je complète le programme}

\textbf{Banque de mots :} input, int, print, prenom, age.

\begin{python}[niveauA.py]
prenom = _______("Ton prenom : ")
age = _______(input("Ton age : "))
_______(f"{prenom} a {age} ans.")
\end{python}

\subsection*{Activité : relie chaque valeur à son type}
\begin{tabularx}{\linewidth}{Y Y}
\toprule
\textbf{Valeur} & \textbf{Type (int, float, str ou bool)}\\
\midrule
\texttt{42} & \\
\texttt{"Bonjour"} & \\
\texttt{3.5} & \\
\texttt{True} & \\
\texttt{"12"} & \\
\bottomrule
\end{tabularx}

\subsection*{Activité : prédis puis teste}
\texttt{print(2 + 3)} affiche \dotfill\quad \texttt{print("2" + "3")} affiche \dotfill

%==================================================================
\newpage
\section*{Annexe D : Niveau B (au niveau)}
\entete{blue}{Niveau B : la moyenne de trois notes}

\textbf{Cahier des charges :} le programme demande trois notes, calcule la moyenne, l'arrondit à une décimale avec \texttt{round(valeur, 1)} et affiche une phrase.

\begin{tabularx}{\linewidth}{>{\bfseries}p{3cm}Y Y Y}
\toprule
Test & Notes saisies & Résultat attendu & Résultat obtenu\\
\midrule
1 & 10, 12, 14 & & \\
2 & 15, 15, 15 & & \\
3 & 8, 9, 12 & & \\
\bottomrule
\end{tabularx}

\subsection*{Questions}
\begin{enumerate}
 \item Pourquoi faut-il convertir les notes avec \texttt{float} ou \texttt{int} ? \dotfill
 \item Que se passe-t-il si tu tapes une lettre au lieu d'un nombre ? \dotfill
\end{enumerate}

%==================================================================
\newpage
\section*{Annexe E : Niveau C (au-dessus du niveau)}
\entete{violet}{Niveau C : chasse aux bogues et convertisseur}

\subsection*{Mission 1 : trouve et corrige trois bogues}
\begin{python}[niveauC.py]
nom = input("Ton nom : ")
nombre = input("Un nombre : ")
double = nombre * 2
print(f"{Nom} : le double de {nombre} est {double}")
\end{python}

\begin{tabularx}{\linewidth}{>{\bfseries}p{1.6cm}Y Y}
\toprule
Bogue & Ce qui ne va pas & Ma correction\\
\midrule
1 & & \\ 2 & & \\ 3 & & \\
\bottomrule
\end{tabularx}

\subsection*{Mission 2 : convertisseur Celsius vers Fahrenheit}
Formule : $F = C \times 1{,}8 + 32$. Écris le programme avec saisie et affichage d'une phrase.\\
\lig \lig \lig

\textbf{Défi bonus :} ajoute la conversion inverse et un affichage avec une seule décimale.

%==================================================================
\newpage
\section*{Annexe F : Corrections et attendus (enseignant)}

\subsection*{Annexe B}
3a : \texttt{int}, \texttt{prenom}. 3b : \texttt{longueur = float(input(...))}, \texttt{largeur = float(input(...))}, \texttt{aire = longueur * largeur}, \texttt{perimetre = 2 * (longueur + largeur)}. Types : float, float, float, float.

\subsection*{Annexe C}
\texttt{input}, \texttt{int}, \texttt{print}. Types : \texttt{42} int ; \texttt{"Bonjour"} str ; \texttt{3.5} float ; \texttt{True} bool ; \texttt{"12"} str. \texttt{print(2 + 3)} affiche 5 ; \texttt{print("2" + "3")} affiche 23.

\subsection*{Annexe D}
\texttt{moyenne = round((n1 + n2 + n3) / 3, 1)}. Résultats : 12,0 ; 15,0 ; 9,7. Sans conversion, \texttt{input} renvoie du texte et les notes se concatènent. Une lettre provoque une \texttt{ValueError} (la gestion d'erreur sera vue plus tard).

\subsection*{Annexe E}
Bogues : 1 \texttt{nombre} est du texte : \texttt{int(input(...))} ; 2 \texttt{double = nombre * 2} multiplierait le texte, corrigé par la conversion ; 3 \texttt{Nom} avec majuscule au lieu de \texttt{nom}. Convertisseur : \texttt{f = c * 1.8 + 32}.\\
\textbf{Critères de réussite :} trois bogues expliqués, convertisseur correct et testé (0 \degre{}C donne 32 \degre{}F).

%==================================================================
\newpage
\section*{Annexe G : Exit ticket et grille d'observation}
\entete{gray}{Exit ticket : variables et types}

\begin{enumerate}
 \item Quel est le type de la variable \texttt{x} après \texttt{x = input("Nombre : ")} ?\\ \lig
 \item Que vaut \texttt{int("4") + 3} ? Et \texttt{"4" + "3"} ?\\ \lig
 \item \autoeval
\end{enumerate}
\textbf{Attendus :} 1. \texttt{str} (du texte). 2. 7, puis \texttt{"43"}.

\bigskip
\grille{Type identifié}{Conversion correcte}{Programme testé}

\end{document}
